\section{GUI Development(Mert Ali Turk and Mohammad Mahdi Mohammadi)}

\subsection{Introduction to GUI}
\label{subsection:intro}

The Measurement Network Gateway (MNG) Client GUI is a real-time sensor
data visualization and control application developed in the C programming
language, utilizing the GTK3 toolkit for the user interface and the Cairo
2D graphics library for real-time chart rendering. The application serves
as the operator-facing front-end of the broader MNG system, establishing
a persistent TCP/IP connection to an MNG server and enabling the user to
monitor, configure, and control up to five sensor channels ---
Temperature, ADC0, ADC1, Switch, and Push Button --- from a single
integrated window.

A fundamental design principle of the application is the strict
separation between the background network thread and the GTK main thread.
All widget updates are dispatched exclusively from the GTK main thread,
ensuring full GTK thread safety without blocking the user interface during
data reception. Access to the shared graph data buffers is synchronized
using a dedicated \texttt{pthread\_mutex\_t}, which serializes writes from
the network thread against reads performed by the Cairo draw callback.
Communication with the MNG server follows a binary protocol over a
persistent TCP connection on port 50012, supporting commands for data
requests, per-sensor frequency configuration, and remote server shutdown.


\subsection{Implementation of GUI}
\subsubsection{GUI Layout Implementation}

The graphical user interface of the MNG Client is organized into three
main sections arranged vertically within the main window. The layout
follows a hierarchical GTK3 widget tree, with a vertical box
(\texttt{GtkBox}) as the root container into which all sections are
packed sequentially from top to bottom.

At the top of the window, a \textit{Connection} frame provides the
controls necessary to establish and manage communication with the server,
including server address configuration, connection management, and
data acquisition control. A status label beneath this frame reflects the
current connection state in real time.

Below the connection controls, a \textit{Sensor Frequency Set} frame
allows the operator to individually configure the sampling rate of each
sensor channel. Each channel accepts a frequency value in the range of
10 to 10{,}000\,Hz with a default of 500\,Hz. The operator can also
adjust the sliding time window of the graph, with a default duration of
5{,}000{,}000\,\textmu s (5 seconds), and issue a remote shutdown command
to the server from this section.

The lower portion of the window is divided horizontally into two panels.
The left panel displays the most recently received numerical values and
relative timestamps for each sensor channel, including Temperature, ADC0,
ADC1, Switch, and Push Button. The right panel contains the Cairo drawing
area for real-time graph rendering, where the operator can select up to
two sensor channels for simultaneous plotting and control the graph
display.

The following code shows the declaration of the widgets used in the GUI, which are organized into the three main sections described above.
\begin{minted}[bgcolor = lightgray, tabsize = 4, breaklines]{c}
typedef struct {
    GtkWidget *window;
    GtkWidget *drawing_area;

    /* Sensor labels */
    GtkWidget *label_temp;
    GtkWidget *label_temp_ts;
    GtkWidget *label_adc0;
    GtkWidget *label_adc0_ts;
    GtkWidget *label_adc1;
    GtkWidget *label_adc1_ts;
    GtkWidget *label_switch;
    GtkWidget *label_pushb;
    GtkWidget *label_status;

    /* Connection controls */
    GtkWidget *entry_ip;
    GtkWidget *button_connect;
    GtkWidget *button_disconnect;
    GtkWidget *button_start;
    GtkWidget *button_stop;

    /* Configure controls */
    GtkWidget *spin_temp_cfg;
    GtkWidget *spin_adc0_cfg;
    GtkWidget *spin_adc1_cfg;
    GtkWidget *spin_switch_cfg;
    GtkWidget *spin_pushb_cfg;
    GtkWidget *button_configure;
    GtkWidget *spin_sliding_window;

    /* Graph controls */
    GtkWidget *combo_sensor1;
    GtkWidget *combo_sensor2;
    GtkWidget *button_graph;
    GtkWidget *button_clear;
    GtkWidget *label_graph_status;

    GtkWidget *btn_temp_set;
    GtkWidget *btn_adc0_set;
    GtkWidget *btn_adc1_set;
    GtkWidget *btn_switch_set;
    GtkWidget *button_shutdown;
    GtkWidget *btn_sliding_window_set;
    GtkWidget *toggle_sampling;

    GtkWidget *spin_req_data_time;
    GtkWidget *btn_req_data_time_set;


    int socket_fd;
    pthread_t receive_thread;
    int connected;
    int thread_running;
    int transmitting;

    sensor_graph_data graph_data[2];
    pthread_mutex_t graph_mutex;

    int current_sensor1;
    int current_sensor2;

    /* Cache last displayed values to prevent redundant updates */
    uint32_t display_temp;
    uint64_t display_temp_ts;
    uint32_t display_adc0;
    uint64_t display_adc0_ts;
    uint32_t display_adc1;
    uint64_t display_adc1_ts;
    uint32_t display_swbtn;
    uint32_t display_pushb;

    guint timer_id;
    uint32_t interval_ms;

    uint64_t start_time;
    pthread_mutex_t display_mutex;
} gui_context;
\end{minted}

\subsubsection{Callback Implmentation}
\subsubsection{Connection button callback}
The \texttt{on\_connect\_clicked()} function is invoked when the operator
presses the \textit{Connect} button and is responsible for establishing
the TCP connection between the GUI client and the MNG server.

The function begins by validating the current application state and user
input. If a connection is already active, the function returns immediately
without taking any action. Otherwise, it reads the server IP address from
the input field and verifies that it is not empty before proceeding.

A TCP socket is then created using the \texttt{socket()} system call with
the \texttt{AF\_INET} and \texttt{SOCK\_STREAM} parameters, specifying an
IPv4 TCP connection. The server address structure is populated with the
target IP address, converted from text to binary form using
\texttt{inet\_pton()}, and the fixed server port defined by
\texttt{SERVER\_PORT}. The \texttt{connect()} system call is then issued
to establish the connection. If any of these steps fail, an appropriate
error message is displayed in the status label, the socket is closed, and
the function returns without modifying the application state.

Upon a successful connection, the application context flags
\texttt{connected} and \texttt{thread\_running} are set, and the
background receive thread is launched via \texttt{pthread\_create()}.
This thread runs independently of the GTK main thread and is responsible
for receiving incoming sensor data from the server. If thread creation
fails, the connection is torn down and the error is reported to the
operator.

Finally, the function updates the user interface to reflect the connected
state by enabling the controls that require an active connection, such as
the data acquisition, configuration, and shutdown buttons, while disabling
the IP entry field and the \textit{Connect} button to prevent duplicate
connection attempts.

\begin{minted}[bgcolor = lightgray, tabsize = 4, breaklines]{c}
void on_connect_clicked(GtkButton *button, gpointer data)
{
    if (ctx.connected) {
        gtk_label_set_text(GTK_LABEL(ctx.label_status), "Already connected!");
        return;
    }

    const char *server_ip = gtk_entry_get_text(GTK_ENTRY(ctx.entry_ip));
    if (!server_ip || strlen(server_ip) == 0) {
        gtk_label_set_text(GTK_LABEL(ctx.label_status), "Enter server IP address!");
        return;
    }

    log_debug("Connecting to %s:%d", server_ip, SERVER_PORT);

    ctx.socket_fd = socket(AF_INET, SOCK_STREAM, 0);
    if (ctx.socket_fd < 0) {
        perror("[GUI] Socket creation error");
        gtk_label_set_text(GTK_LABEL(ctx.label_status), "Socket error!");
        return;
    }

    struct sockaddr_in serv_addr;
    memset(&serv_addr, 0, sizeof(serv_addr));
    serv_addr.sin_family = AF_INET;
    serv_addr.sin_port = htons(SERVER_PORT);

    if (inet_pton(AF_INET, server_ip, &serv_addr.sin_addr) <= 0) {
        perror("[GUI] Invalid address");
        gtk_label_set_text(GTK_LABEL(ctx.label_status), "Invalid server address!");
        close(ctx.socket_fd);
        ctx.socket_fd = -1;
        return;
    }

    if (connect(ctx.socket_fd, (struct sockaddr *)&serv_addr, sizeof(serv_addr)) < 0) {
        perror("[GUI] Connection error");
        gtk_label_set_text(GTK_LABEL(ctx.label_status), "Connection failed!");
        close(ctx.socket_fd);
        ctx.socket_fd = -1;
        return;
    }

    ctx.connected = 1;
    ctx.thread_running = 1;
    ctx.start_time = 0;

    if (pthread_create(&ctx.receive_thread, NULL, receive_data_thread, NULL) < 0) {
        perror("[GUI] Thread creation error");
        gtk_label_set_text(GTK_LABEL(ctx.label_status), "Thread creation error!");
        close(ctx.socket_fd);
        ctx.socket_fd = -1;
        ctx.connected = 0;
        return;
    }

    char status_msg[128];
    snprintf(status_msg, sizeof(status_msg), "Connected to %s!", server_ip);
    gtk_label_set_text(GTK_LABEL(ctx.label_status), status_msg);
    gtk_widget_set_sensitive(GTK_WIDGET(ctx.entry_ip), FALSE);
    gtk_widget_set_sensitive(GTK_WIDGET(ctx.button_connect), FALSE);
    gtk_widget_set_sensitive(GTK_WIDGET(ctx.button_disconnect), TRUE);
    gtk_widget_set_sensitive(GTK_WIDGET(ctx.button_start), TRUE);
    gtk_widget_set_sensitive(GTK_WIDGET(ctx.button_configure), TRUE);
    gtk_widget_set_sensitive(GTK_WIDGET(ctx.btn_temp_set), TRUE);
    gtk_widget_set_sensitive(GTK_WIDGET(ctx.btn_adc0_set), TRUE);
    gtk_widget_set_sensitive(GTK_WIDGET(ctx.btn_adc1_set), TRUE);
    gtk_widget_set_sensitive(GTK_WIDGET(ctx.btn_switch_set), TRUE);
    gtk_widget_set_sensitive(GTK_WIDGET(ctx.button_shutdown), TRUE);
    gtk_widget_set_sensitive(ctx.btn_req_data_time_set, TRUE);
    gtk_widget_set_sensitive(ctx.spin_req_data_time, TRUE);
    gtk_widget_set_sensitive(ctx.btn_req_data_time_set, TRUE);
    gtk_widget_set_sensitive(ctx.spin_req_data_time, TRUE);



    log_debug("Connected successfully");
}
\end{minted}

\subsubsection{Disconnection Callback}

The \texttt{on\_disconnect\_clicked()} function is invoked when the
operator presses the \textit{Disconnect} button and is responsible for
cleanly terminating the TCP connection with the MNG server and restoring
the application to its idle state. If no active connection exists, the
function returns immediately without taking any action.

The disconnection sequence begins by clearing the \texttt{thread\_running}
and \texttt{transmitting} flags to signal the background receive thread
to stop. A \texttt{disconnect} command is sent to the server, followed by
a short delay of 100 milliseconds to allow pending socket data to be
flushed before the socket is closed. \texttt{pthread\_join()} is then
called to wait for the receive thread to fully terminate, ensuring that
no dangling thread accesses the application context after disconnection.

Finally, the user interface is updated to reflect the disconnected state.
Controls that require an active connection are disabled, the IP entry
field and the \textit{Connect} button are re-enabled, and the
\textit{Start} button label is reset to its default text to prepare the
interface for a new connection session.

\begin{minted}[bgcolor = lightgray, tabsize = 4, breaklines]{c}
void on_disconnect_clicked(GtkButton *button, gpointer data)
{
    if (!ctx.connected) {
        gtk_label_set_text(GTK_LABEL(ctx.label_status), "Not connected!");
        return;
    }

    log_debug("Disconnecting...");

    ctx.thread_running = 0;
    ctx.transmitting = 0;

    const char *disconnect_cmd = "disconnect\n";
    if (send(ctx.socket_fd, disconnect_cmd, strlen(disconnect_cmd), 0) < 0) {
        log_error("Failed to send disconnect command");
    }

    usleep(100000);

    if (ctx.socket_fd >= 0) {
        close(ctx.socket_fd);
        ctx.socket_fd = -1;
    }

    pthread_join(ctx.receive_thread, NULL);

    ctx.connected = 0;

    gtk_label_set_text(GTK_LABEL(ctx.label_status), "Disconnected!");
    gtk_widget_set_sensitive(GTK_WIDGET(ctx.entry_ip), TRUE);
    gtk_widget_set_sensitive(GTK_WIDGET(ctx.button_connect), TRUE);
    gtk_widget_set_sensitive(GTK_WIDGET(ctx.button_disconnect), FALSE);
    gtk_widget_set_sensitive(GTK_WIDGET(ctx.button_start), FALSE);
    gtk_widget_set_sensitive(GTK_WIDGET(ctx.button_stop), FALSE);
    gtk_widget_set_sensitive(GTK_WIDGET(ctx.button_configure), FALSE);
    gtk_widget_set_sensitive(GTK_WIDGET(ctx.btn_temp_set), FALSE);
    gtk_widget_set_sensitive(GTK_WIDGET(ctx.btn_adc0_set), FALSE);
    gtk_widget_set_sensitive(GTK_WIDGET(ctx.btn_adc1_set), FALSE);
    gtk_widget_set_sensitive(GTK_WIDGET(ctx.btn_switch_set), FALSE);
    gtk_widget_set_sensitive(ctx.btn_req_data_time_set, FALSE);
    gtk_widget_set_sensitive(ctx.spin_req_data_time, FALSE);


    gtk_button_set_label(GTK_BUTTON(ctx.button_start), "Start");

    log_debug("Disconnected successfully");
}
\end{minted}

\subsubsection{Start Callback}

The \texttt{on\_start\_clicked()} function initiates the sensor data
acquisition process when the operator presses the \textit{Start} button.
It first verifies that a connection is active and that no transmission is
already in progress, returning immediately if either condition is not met.

An initial data request command \texttt{MNG\_CMD\_REQ\_DATA} is
transmitted to the server using the binary \texttt{Clnt\_Send\_struct}
protocol structure. Both graph data buffers are cleared beforehand to
remove any stale data from a previous session. Upon success, the
\texttt{transmitting} flag is set and a periodic GTK timer is started via
\texttt{g\_timeout\_add()}, which repeatedly calls
\texttt{data\_request\_timer()} at a fixed interval to implement the
request-response polling mechanism. The user interface is then updated
to reflect the active acquisition state.


\begin{minted}[bgcolor = lightgray, tabsize = 4, breaklines]{c}
void on_start_clicked(GtkButton *button, gpointer data)
{
    if (!ctx.connected) {
        gtk_label_set_text(GTK_LABEL(ctx.label_status), "Not connected!");
        return;
    }

    if (ctx.transmitting) {
        gtk_label_set_text(GTK_LABEL(ctx.label_status), "Already transmitting!");
        return;
    }


	Clnt_Send_struct sendCmd;
    sendCmd.cmd = MNG_CMD_REQ_DATA;
    sendCmd.param[0] = 0;
    sendCmd.param[1] = 0;

    if (send(ctx.socket_fd, &sendCmd, sizeof(sendCmd), 0) < 0)
    {
        perror("[GUI] Error sending start command");
        gtk_label_set_text(GTK_LABEL(ctx.label_status), "Error sending command!");
        return;
    }

    memset(ctx.graph_data[0].points, 0, sizeof(ctx.graph_data[0].points));
    memset(ctx.graph_data[1].points, 0, sizeof(ctx.graph_data[1].points));

    ctx.transmitting = 1;
    ctx.timer_id = g_timeout_add(INITIAL_TIME_INTERVAL_REQ, (GSourceFunc)data_request_timer, NULL);


    gtk_label_set_text(GTK_LABEL(ctx.label_status), "Data acquisition started!");
    gtk_widget_set_sensitive(GTK_WIDGET(ctx.button_start), FALSE);
    gtk_widget_set_sensitive(GTK_WIDGET(ctx.button_stop), TRUE);
    gtk_button_set_label(GTK_BUTTON(ctx.button_start), "Started");


    log_debug("Transmission started");
}
\end{minted}

\subsubsection{Stop Callback}

The \texttt{on\_stop\_clicked()} function halts the ongoing data
acquisition when the operator presses the \textit{Stop} button. If no
active transmission is in progress, the function returns immediately
without taking any action.

The \texttt{transmitting} flag is cleared, which causes the periodic
polling timer to stop at its next expiry, as \texttt{data\_request\_timer()}
checks this flag before sending each request. The user interface is then
updated by re-enabling the \textit{Start} button and disabling the
\textit{Stop} button, restoring the interface to its connected idle state.


\begin{minted}[bgcolor = lightgray, tabsize = 4, breaklines]{c}
void on_stop_clicked(GtkButton *button, gpointer data)
{
    if (!ctx.connected || !ctx.transmitting) {
        gtk_label_set_text(GTK_LABEL(ctx.label_status), "Not transmitting!");
        return;
    }
    
    ctx.transmitting = 0;
    gtk_label_set_text(GTK_LABEL(ctx.label_status), "Data acquisition stopped!");
    gtk_widget_set_sensitive(GTK_WIDGET(ctx.button_start), TRUE);
    gtk_widget_set_sensitive(GTK_WIDGET(ctx.button_stop), FALSE);
    gtk_button_set_label(GTK_BUTTON(ctx.button_start), "Start");

    log_debug("Transmission stopped");
}
\end{minted}

\subsubsection{Sensor Frequency Configuration Callback}

Each sensor channel has a dedicated frequency configuration callback,
all following the same structure. The \texttt{on\_set\_temp\_clicked()}
function is presented here as a representative example.

The function reads the desired sampling frequency from the corresponding
spin button. A value of zero sends a dormant command to the server,
disabling sampling for that channel. For a non-zero value, the function
validates that the frequency evenly divides 1000\,Hz before constructing
and transmitting a \texttt{MNG\_CMD\_SET\_FREQ} command to the server
using the \texttt{Clnt\_Send\_struct} protocol structure. The status
label is updated to confirm the outcome.

\begin{minted}[bgcolor = lightgray, tabsize = 4, breaklines]{c}
void on_set_temp_clicked(GtkButton *button, gpointer data)
{
    if (!ctx.connected)
    {
        gtk_label_set_text(GTK_LABEL(ctx.label_status), "Not connected!");
        return;
    }

    Clnt_Send_struct sendCmd;
    int temp_cfg = (int)gtk_spin_button_get_value(GTK_SPIN_BUTTON(ctx.spin_temp_cfg));


    if(temp_cfg == 0)
    {
    	sendCmd.cmd = MNG_CMD_DORMANT_TEM;
        sendCmd.param[0] = 0;
        sendCmd.param[1] = 0;

    }else
    {
        if ((1000 % temp_cfg) != 0) {
            gtk_label_set_text(GTK_LABEL(ctx.label_status),
                               "Invalid TEMP frequency: must divide 1000 Hz");
            return;
        }

        sendCmd.cmd = MNG_CMD_SET_FREQ;
        sendCmd.param[0] = MNG_CMD_SET_FREQ_TMP;
        sendCmd.param[1] = temp_cfg;
    }


    if (send(ctx.socket_fd , &sendCmd, sizeof(Clnt_Send_struct),0) < 0)
    {
        perror("[GUI] Error sending configure command");
        gtk_label_set_text(GTK_LABEL(ctx.label_status), "Error sending configure command!");
        return;
    }

    char status_msg[256];
    snprintf(status_msg, sizeof(status_msg),
             "Configure sent: Temp=%d Hz ",
			 temp_cfg);
    gtk_label_set_text(GTK_LABEL(ctx.label_status), status_msg);
}
\end{minted}


\subsubsection{Graph Plot Callback}

The \texttt{on\_graph\_clicked()} function is invoked when the operator
presses the \textit{Plot Graph} button. It reads the selected sensor
channels from the two combo boxes and validates the input, ensuring that
two different sensors are selected and that both belong to the plottable
channels, which are limited to Temperature, ADC0, and ADC1.

Upon successful validation, the selected sensor identifiers are stored in
the application context and both graph data buffers are reset under mutex
protection to clear any previously accumulated data. A redraw of the
Cairo drawing area is then triggered via \texttt{gtk\_widget\_queue\_draw()},
initiating the rendering of the newly selected sensor pair.


\begin{minted}[bgcolor = lightgray, tabsize = 4, breaklines]{c}
void on_graph_clicked(GtkButton *button, gpointer data)
{
    int sensor1_index = gtk_combo_box_get_active(GTK_COMBO_BOX(ctx.combo_sensor1));
    int sensor2_index = gtk_combo_box_get_active(GTK_COMBO_BOX(ctx.combo_sensor2));

    if (sensor1_index < 0 || sensor2_index < 0) {
        gtk_label_set_text(GTK_LABEL(ctx.label_graph_status), "Select two sensors!");
        return;
    }

    if (sensor1_index == sensor2_index) {
        gtk_label_set_text(GTK_LABEL(ctx.label_graph_status), "Select TWO DIFFERENT sensors!");
        return;
    }

    int sensor1 = sensor1_index + 1;
    int sensor2 = sensor2_index + 1;

    if (sensor1 < SENSOR_TEMP || sensor1 > SENSOR_ADC1 ||
        sensor2 < SENSOR_TEMP || sensor2 > SENSOR_ADC1) {
        gtk_label_set_text(GTK_LABEL(ctx.label_graph_status), "Only Temperature, ADC0, ADC1 can be graphed!");
        return;
    }

    ctx.current_sensor1 = sensor1;
    ctx.current_sensor2 = sensor2;

    pthread_mutex_lock(&ctx.graph_mutex);
    ctx.graph_data[0].count = 0;
    ctx.graph_data[0].min_val = UINT32_MAX;
    ctx.graph_data[0].max_val = 0;
    ctx.graph_data[1].count = 0;
    ctx.graph_data[1].min_val = UINT32_MAX;
    ctx.graph_data[1].max_val = 0;
    pthread_mutex_unlock(&ctx.graph_mutex);

    char status_msg[256];
    snprintf(status_msg, sizeof(status_msg),
             "Plotting %s (red) vs %s (blue)...",
             sensor_names[sensor1],
             sensor_names[sensor2]);
    gtk_label_set_text(GTK_LABEL(ctx.label_graph_status), status_msg);
    gtk_widget_queue_draw(ctx.drawing_area);

    log_debug("Graph: %s vs %s", sensor_names[sensor1], sensor_names[sensor2]);
}
\end{minted}


\subsubsection{Clear Graph Callback}

The \texttt{on\_clear\_graph\_clicked()} function resets the graph display
when the operator presses the \textit{Clear Data} button. Both graph data
buffers are fully cleared under mutex protection, with their sample
counts, minimum and maximum values, and data point arrays all reset to
their initial state. A redraw of the Cairo drawing area is then triggered
to immediately reflect the cleared state on the screen.

\begin{minted}[bgcolor = lightgray, tabsize = 4, breaklines]{c}
void on_clear_graph_clicked(GtkButton *button, gpointer data)
{
    pthread_mutex_lock(&ctx.graph_mutex);
    ctx.graph_data[0].count = 0;
    ctx.graph_data[0].min_val = UINT32_MAX;
    ctx.graph_data[0].max_val = 0;
    ctx.graph_data[1].count = 0;
    ctx.graph_data[1].min_val = UINT32_MAX;
    ctx.graph_data[1].max_val = 0;
    memset(ctx.graph_data[0].points, 0, sizeof(ctx.graph_data[0].points));
    memset(ctx.graph_data[1].points, 0, sizeof(ctx.graph_data[1].points));
    pthread_mutex_unlock(&ctx.graph_mutex);

    gtk_label_set_text(GTK_LABEL(ctx.label_graph_status), "Graph data cleared!");
    gtk_widget_queue_draw(ctx.drawing_area);

    log_debug("Graph cleared");
}
\end{minted}


\subsubsection{Shutdown Callback}

The \texttt{on\_shutdown\_clicked()} function transmits a remote shutdown
command to the MNG server when the operator presses the \textit{Shutdown}
button. If no active connection exists, the function returns immediately
without taking any action.

A \texttt{MNG\_CMD\_SHUTDOWN} command is constructed using the
\texttt{Clnt\_Send\_struct} protocol structure and transmitted to the
server over the active TCP socket. Upon success, the connection,
disconnection, and start buttons are disabled, reflecting the fact that
the server process is no longer available to accept further commands.


\begin{minted}[bgcolor = lightgray, tabsize = 4, breaklines]{c}
void on_shutdown_clicked(GtkWidget *widget, gpointer data)
{
    if (!ctx.connected)
    {
        gtk_label_set_text(GTK_LABEL(ctx.label_status), "Not connected!");
        return;
    }

    Clnt_Send_struct sendCmd;
    sendCmd.cmd = MNG_CMD_SHUTDOWN;
    sendCmd.param[0] = 0;
    sendCmd.param[1] = 0;

    if (send(ctx.socket_fd , &sendCmd, sizeof(Clnt_Send_struct),0) < 0)
    {
        perror("[GUI] Error sending configure command");
        gtk_label_set_text(GTK_LABEL(ctx.label_status), "Error sending configure command!");
        return;
    }

    gtk_widget_set_sensitive(GTK_WIDGET(ctx.button_connect), FALSE);
    gtk_widget_set_sensitive(GTK_WIDGET(ctx.button_disconnect), FALSE);
    gtk_widget_set_sensitive(GTK_WIDGET(ctx.button_start), FALSE);

}
\end{minted}


\subsubsection{Sliding Window Callback}

The \texttt{on\_set\_sliding\_window\_clicked()} function updates the
graph's sliding time window duration. The requested value is read from
the spin button and validated to be within the accepted range of
100{,}000 to 10{,}000{,}000\,\textmu s. If valid, the global
\texttt{SLIDING\_WINDOW\_US} variable is updated and confirmed in the
status label, otherwise an error message is displayed and the value
remains unchanged.
\begin{minted}[bgcolor = lightgray, tabsize = 4, breaklines]{c}
void on_set_sliding_window_clicked(GtkButton *button, gpointer data)
{
    uint64_t new_window;


    new_window = (uint64_t)gtk_spin_button_get_value(GTK_SPIN_BUTTON(ctx.spin_sliding_window));



    if (new_window < 100000 || new_window > 10000000) {
        gtk_label_set_text(GTK_LABEL(ctx.label_status),
                           "Sliding window must be 100,000 – 10,000,000 µs");
        return;
    }

    SLIDING_WINDOW_US = new_window;

    char msg[128];
    snprintf(msg, sizeof(msg), "Sliding window set to %" PRIu64 " µs", SLIDING_WINDOW_US);
    gtk_label_set_text(GTK_LABEL(ctx.label_status), msg);
}
\end{minted}

\subsubsection{Sensor Selection Callbacks}

Two sensor selection callbacks, \texttt{on\_sensor1\_changed()} and
\texttt{on\_sensor2\_changed()}, are provided to handle the selection of
the two sensor channels that can be simultaneously plotted on the graph,
corresponding to the red and blue plot lines respectively. When the
operator changes the selection in either combo box, the selected index is
mapped to the corresponding sensor identifier and stored in the
application context. The associated graph data buffer is immediately
cleared and a redraw is triggered to reflect the updated selection on
the graph.


\begin{minted}[bgcolor = lightgray, tabsize = 4, breaklines]{c}
void on_sensor1_changed(GtkComboBox *combo, gpointer data)
{
    int idx = gtk_combo_box_get_active(combo);

    switch (idx)
    {
		case 0:
			ctx.current_sensor1 = SENSOR_NONE;
			break;
		case 1:
			ctx.current_sensor1 = SENSOR_TEMP;
			break;
		case 2:
			ctx.current_sensor1 = SENSOR_ADC0;
			break;
		case 3:
			ctx.current_sensor1 = SENSOR_ADC1;
			break;
		default:
			return;
    }

    memset(ctx.graph_data[0].points, 0, sizeof(ctx.graph_data[0].points));
    ctx.graph_data[0].count = 0;

    gtk_label_set_text(GTK_LABEL(ctx.label_graph_status),
                       "Sensor 1 updated");

    gtk_widget_queue_draw(ctx.drawing_area);
}
\end{minted}

\subsubsection{Data Request Interval Callback}

The \texttt{on\_set\_req\_data\_time\_clicked()} function allows the
operator to adjust the polling interval at runtime. The desired interval
is read from the corresponding spin button in milliseconds. If a polling
timer is already active, it is removed before a new timer is registered
via \texttt{g\_timeout\_add()} with the updated interval. This allows the
operator to dynamically control the rate at which data requests are sent
to the server without interrupting the ongoing acquisition.

\begin{minted}[bgcolor = lightgray, tabsize = 4, breaklines]{c}
void on_set_req_data_time_clicked(GtkButton *button, gpointer user_data)
{
    int req_time_ms =
        gtk_spin_button_get_value_as_int(
            GTK_SPIN_BUTTON(ctx.spin_req_data_time));

    if (ctx.timer_id > 0)
    {
        g_source_remove(ctx.timer_id);
        ctx.timer_id = 0;
    }

    ctx.interval_ms = req_time_ms;
    ctx.timer_id = g_timeout_add(req_time_ms, (GSourceFunc)data_request_timer, NULL);
    log_debug("Data request timer updated to %u ms", req_time_ms);
}
\end{minted}


\subsubsection{Window Destroy Callback}

The \texttt{on\_window\_destroy()} function is invoked when the operator
closes the main application window. If an active connection exists, it
calls \texttt{on\_disconnect\_clicked()} to ensure the TCP connection is
cleanly terminated and the background receive thread is stopped before
the application exits. The GTK main loop is then terminated via
\texttt{gtk\_main\_quit()}.

\begin{minted}[bgcolor = lightgray, tabsize = 4, breaklines]{c}
void on_window_destroy(GtkWidget *widget, gpointer data)
{
    if (ctx.connected) {
        on_disconnect_clicked(NULL, NULL);
    }
    gtk_main_quit();
}
\end{minted}



\subsubsection{Network Thread}
The network thread is responsible for establishing and maintaining the TCP/IP connection with the MNG server, receiving sensor data, and updating the GUI accordingly. It operates in a continuous loop, blocking on socket reads until data is received. Upon receiving a data packet, it parses the sensor information based on predefined IDs and updates the corresponding GUI elements with the new sensor values and timestamps. The thread also ensures that any updates to the graph data structure are performed in a thread-safe manner by acquiring the appropriate mutex locks before modifying shared resources. This design allows for real-time data visualization without blocking the main GTK thread, ensuring a responsive user interface.

\begin{minted}[bgcolor = lightgray, tabsize = 4, breaklines]{c}
void *receive_data_thread(void *arg)
{
    log_debug("Receive thread started");

   // unsigned char raw_payload[PAYLOAD_SIZE];
    mng_Payload_t response;



    while (ctx.thread_running) {
        //memset(raw_payload, 0, PAYLOAD_SIZE);

        log_debug("Waiting for payload...");
        memset(&response, 0, sizeof(response));
        ssize_t bytes_received = recv(ctx.socket_fd, &response, sizeof(response), 0);

        log_debug("Received %zd bytes", bytes_received);

        if (bytes_received < 0) {
            if (errno != ECONNRESET && errno != EPIPE) {
                perror("[GUI] recv error");
            }
            log_debug("recv() returned error, breaking");
            break;
        } else if (bytes_received == 0) {
            log_debug("Server closed connection");
            break;
        }

       /* if (bytes_received != PAYLOAD_SIZE) {
            log_debug("WARNING: Received %zd bytes, expected %d", bytes_received, PAYLOAD_SIZE);
        } */

        int packets_parsed = 0;
        log_debug("Total getting sample count = %u",  response.sampleCount);

        for (uint32_t i = 0; i < response.sampleCount  ; i++)
        {

            //packet *pkt = (packet *)&raw_payload[i * PACKET_SIZE];

//            log_debug("Packet %d: sid=%u, svalue=%u, ts=%" PRIu64,
//                     i, response.sendBuf[i].sensorID,
//					 	response.sendBuf[i].sensorVal,
//						response.sendBuf[i].elapsedTime);

            if ( response.sendBuf[i].sensorID > 5) {
                log_error("Invalid sensor ID in packet %d: %u", i, response.sendBuf[i].sensorID);
                continue;
            }

            if (ctx.start_time == 0) {
                ctx.start_time = response.sendBuf[i].elapsedTime;
                log_debug("Start time set to: %" PRIu64, ctx.start_time);
            }

            sensor_data_update *update = (sensor_data_update *)g_malloc(sizeof(sensor_data_update));
            if (update == NULL) {
                log_error("Failed to allocate memory for sensor_data_update");
                break;
            }

            update->sid    = response.sendBuf[i].sensorID;
            update->svalue = response.sendBuf[i].sensorVal;
            update->ts     = response.sendBuf[i].elapsedTime;

            log_debug("Queuing update for sensor %d", response.sendBuf[i].sensorID);

            guint idle_id = g_idle_add(update_sensor_display, (gpointer)update);
            if (idle_id == 0) {
                log_error("g_idle_add() failed");
                g_free(update);
                break;
            }

            packets_parsed++;
        }

        log_debug("Parsed %d packets from payload", packets_parsed);
    }

    log_debug("Receive thread exiting");
    return NULL;
}
\end{minted}

\subsubsection{Graph Drawing Methodology}

This section describes the implementation of the drawing mechanism within the GUI. 
In addition, it presents an overview of the key APIs used for graph rendering.First, the graph drawing 
workflow is illustrated in the following Figure \ref{fig:dataflow}.

\begin{figure}[H]
    \centering
    \includegraphics[width=0.7\linewidth]{flow.png}
    \caption{Dataflow of Drawing}
    \label{fig:dataflow}
\end{figure}

The GUI is designed with simplicity in mind. The workflow illustrated above is briefly 
explained as follows. First, the receive thread blocks until data are received from the 
server, which in this case is the Zynq platform. After a data packet is received, the 
sensor display is updated and the data are parsed based on their sensor IDs. The parsed 
sensor data, along with their timestamps, are then added to the graph data structure. 
Finally, all collected data are rendered in the graph.

In the GUI implementation, one critical aspect is the handling of the x-axis timing. 
The GUI must be able to adjust the time window according to user requirements. 
For example, a user may want to view a narrower or wider time window depending on the 
signal being examined. Another important aspect is displaying the most recently sampled 
sensor data within the selected time window. The time window interval is controlled by a 
global variable that can be adjusted at run-time by GUI.

At most, two signals can be drawn in the GUI. Each signal structure consists of several 
attributes: a data point array, a count value, a sensor ID, and minimum and maximum 
values. In the following subsections, the purpose of these attributes is explained in 
detail, with particular emphasis on the minimum and maximum values, since a normalization 
method is applied on the y-axis.

\subsubsection{Update Sensor Display API}

Each sensor sample is first passed to this API, where essential sample details are 
stored in the GUI data structures. Sensors are parsed based on their IDs, after which 
sample details such as the value, timestamp, and sensor ID are assigned to the 
corresponding GUI structures.

After parsing, the GUI is also capable of displaying sensor values and timestamps as 
raw text. Since only two graphs are available, the user may want to view the values of 
additional sensors in textual form. It is important to note that the number of sensors 
may be higher. Finally, the system checks whether the parsed data correspond to the 
currently selected graphs; if so, the \text{add\_graph\_data\_point()}  API is invoked to add the 
parsed data point to the graph.

The \text{add\_graph\_data\_point()} API is explained in the next subsection. The 
following Listing \ref{lst:update_sensor_display} illustrates the implementation of \text{update\_sensor\_display} API.


\begin{lstlisting}[language=C,
caption={Sensor Update Implementation},
label={lst:update_sensor_display}, ]

gboolean update_sensor_display(gpointer data)
{
    sensor_data_update *update = (sensor_data_update *)data;

    if (update == NULL) {
        log_error("update_sensor_display: NULL pointer");
        return FALSE;
    }

    if (update->sid > 5) {
        log_error("Invalid sensor ID: %d", update->sid);
        g_free(update);
        return FALSE;
    }

    char buffer[256];
    int needs_redraw = 0;

    pthread_mutex_lock(&ctx.display_mutex);

    switch (update->sid) {
		case SENSOR_TEMP: {
			uint64_t relative_ms = (update->ts - ctx.start_time) / 1000;

			ctx.display_temp = update->svalue;
			ctx.display_temp_ts = relative_ms;

			snprintf(buffer, sizeof(buffer), "%u", update->svalue);
			gtk_label_set_text(GTK_LABEL(ctx.label_temp), buffer);

			snprintf(buffer, sizeof(buffer), "%" PRIu64 " ms", relative_ms);
			gtk_label_set_text(GTK_LABEL(ctx.label_temp_ts), buffer);

			log_debug("TEMP: %u @ %" PRIu64 " us (relative: %" PRIu64 " ms)",
					 update->svalue, update->ts, relative_ms);

			if (ctx.current_sensor1 == SENSOR_TEMP)
				add_graph_data_point(SENSOR_TEMP, update->svalue, update->ts, 0);
			if (ctx.current_sensor2 == SENSOR_TEMP)
				add_graph_data_point(SENSOR_TEMP, update->svalue, update->ts, 1);
			needs_redraw = 1;
			break;
		}

		case SENSOR_ADC0: {
			uint64_t relative_ms = (update->ts - ctx.start_time) / 1000;

			ctx.display_adc0 = update->svalue;
			ctx.display_adc0_ts = relative_ms;

			snprintf(buffer, sizeof(buffer), "%u", update->svalue);
			gtk_label_set_text(GTK_LABEL(ctx.label_adc0), buffer);

			snprintf(buffer, sizeof(buffer), "%" PRIu64 " ms", relative_ms);
			gtk_label_set_text(GTK_LABEL(ctx.label_adc0_ts), buffer);

			log_debug("ADC0: %u @ %" PRIu64 " us (relative: %" PRIu64 " ms)",
					 update->svalue, update->ts, relative_ms);


			if (ctx.current_sensor1 == SENSOR_ADC0)
				add_graph_data_point(SENSOR_ADC0, update->svalue, update->ts, 0);
			if (ctx.current_sensor2 == SENSOR_ADC0)
				add_graph_data_point(SENSOR_ADC0, update->svalue, update->ts, 1);
			needs_redraw = 1;
			break;
		}

		case SENSOR_ADC1: {
			uint64_t relative_ms = (update->ts - ctx.start_time) / 1000;


			ctx.display_adc1 = update->svalue;
			ctx.display_adc1_ts = relative_ms;

			snprintf(buffer, sizeof(buffer), "%u", update->svalue);
			gtk_label_set_text(GTK_LABEL(ctx.label_adc1), buffer);

			snprintf(buffer, sizeof(buffer), "%" PRIu64 " ms", relative_ms);
			gtk_label_set_text(GTK_LABEL(ctx.label_adc1_ts), buffer);

			log_debug("ADC1: %u @ %" PRIu64 " us (relative: %" PRIu64 " ms)",
					 update->svalue, update->ts, relative_ms);


			if (ctx.current_sensor1 == SENSOR_ADC1)
				add_graph_data_point(SENSOR_ADC1, update->svalue, update->ts, 0);
			if (ctx.current_sensor2 == SENSOR_ADC1)
				add_graph_data_point(SENSOR_ADC1, update->svalue, update->ts, 1);
			needs_redraw = 1;
			break;
		}

		case SENSOR_SWBTN:
			if (ctx.display_swbtn != update->svalue) {
				ctx.display_swbtn = update->svalue;

				snprintf(buffer, sizeof(buffer), "0x%02X", update->svalue);
				gtk_label_set_text(GTK_LABEL(ctx.label_switch), buffer);
				log_debug("SWITCH: 0x%02X", update->svalue);

				ctx.display_pushb = (update->svalue >> 8);
				snprintf(buffer, sizeof(buffer), "0x%02X",  (update->svalue >> 8));
				gtk_label_set_text(GTK_LABEL(ctx.label_pushb), buffer);
				log_debug("PUSHB: 0x%02X", update->svalue);

			}
			break;

		default:
			log_error("Unknown sensor ID: %d", update->sid);
			break;
    }

    pthread_mutex_unlock(&ctx.display_mutex);

    g_free(update);

    if (needs_redraw && ctx.drawing_area != NULL) {
        gtk_widget_queue_draw(ctx.drawing_area);
    }

    return FALSE;
}

\end{lstlisting}


\subsubsection{Add Sampled Sensor Data to Graph API}

The main responsibility of the \text{add\_graph\_data\_point()} API is to receive sampled data and 
add it to the graph structure. While doing so, it also removes any samples that fall 
outside the current time window interval. Implementing the addition of only those samples 
within the time window is straightforward. First, the API checks whether the timestamp of 
the current sample exceeds the sliding window threshold; if so, it updates the start of 
the time window accordingly. After updating the time window, it scans the buffer for any 
samples that are older than the time window and removes them. Finally, the count variable 
in the graph structure is updated, and the sensor values are stored in the graph structure.

At the end of the API, the minimum and maximum sensor values are updated. This is 
necessary because a normalization method is applied to the y-axis, and these values are 
required for proper scaling. Listing \ref{lst:adddata} presents the complete implementation of the API.


\begin{lstlisting}[language=C,
caption={Add Sampled Sensor Data to Graph Implementation},
label={lst:adddata}, ]
void add_graph_data_point(uint32_t sensor_id, uint32_t raw_value, uint64_t timestamp, int graph_index)
{
	/* Get sensor graph data struct as pointer*/
	sensor_graph_data *g = &ctx.graph_data[graph_index];

    pthread_mutex_lock(&ctx.graph_mutex);
    /* Convert raw to unit data*/
    double value = convert_raw_to_unit(sensor_id, raw_value);


    uint64_t window_start = 0;
    uint32_t write = 0;

    /* Are these points older than our timing window? if older, filter it*/
    if (timestamp > SLIDING_WINDOW_US)
    {
    	/* Is time now greater then window time interval ? */
    	window_start = timestamp - SLIDING_WINDOW_US;

        for (int i = 0; i < g->count; i++)
        {
            if (g->points[i].timestamp >= window_start)
            {
                g->points[write++] = g->points[i];
            }
        }
        /* Our new count value since we added some of them to beginning*/
        g->count = write;
    }

    /* Add new data point to buffer */
    if (g->count < MAX_DATA_POINTS)
    {
         g->points[g->count].timestamp = timestamp;
         g->points[g->count].value = value;
         g->count++;
     }

    /* Update new minimum and maximum values*/
    if (g->count == 1)
    {
        g->min_val  = value;
        g->max_val  = value;
    } else
    {
        if (value < g->min_val)
        					g->min_val = value;
        if (value > g->max_val)
        					g->max_val = value;
    }

    pthread_mutex_unlock(&ctx.graph_mutex);
}
\end{lstlisting}


\subsubsection{Draw Graph API}
The \text{on\_draw\_graph()} API is responsible for rendering the graph in the GUI. It first 
retrieves the size of the drawing area and exits if the dimensions are invalid. To ensure 
thread-safe access, it locks the graph data with a mutex. The background of the graph 
area is filled with white to create a clean canvas. If there is no data in either graph 
buffer, a placeholder message “No data to display” is shown, and the function exits after 
releasing the mutex.

The usable plotting area is calculated by subtracting margins from the widget dimensions. 
The minimum and maximum values of each graph are retrieved and adjusted with a small 
padding to avoid drawing at the edges. One of the most important aspects of this API is 
the handling of timing for the x-axis. The API implements a sliding time window, which 
determines the portion of the data that is visible on the graph. The start of the visible 
time range is calculated based on the timestamp of the latest sampled data. Only data 
points within this time window are drawn, ensuring that the graphs always display the most 
recent and relevant information. This approach allows the GUI to dynamically update in 
real time while providing a consistent view of the signal over a fixed duration.

After setting general plot configurations, the two graphs are drawn using their respective 
data and axes ranges. Labels for the y-axis and x-axis are added to indicate values and 
time intervals, and graph titles are displayed to provide context. Finally, the mutex is 
released, completing the drawing process. By carefully managing the timing and applying 
the sliding window, the GUI ensures that users see up-to-date graph data in a visually 
organized and real-time manner.

Listing \ref{lst:draw} presents the complete implementation of the drawing API.

\begin{lstlisting}[language=C,
caption={Graph Draw API},
label={lst:draw}, ]

gboolean on_draw_graph(GtkWidget *widget, cairo_t *cr, gpointer data)
{
	/* Get Info on size of GUI*/
    int width = gtk_widget_get_allocated_width(widget);
    int height = gtk_widget_get_allocated_height(widget);

    if (width <= 0 || height <= 0)
        					return FALSE;

    pthread_mutex_lock(&ctx.graph_mutex);

    /* Draw the background of GUI */
    cairo_set_source_rgb(cr, 1, 1, 1); /* White color */
    cairo_rectangle(cr, 0, 0, width, height);
    cairo_fill(cr);

    /* If there is no data, write text that no data in buffer! */
    if (ctx.graph_data[0].count == 0 && ctx.graph_data[1].count == 0)
    {
        cairo_set_source_rgb(cr, 0.5, 0.5, 0.5);
        cairo_select_font_face(cr, "Monospace",
                               CAIRO_FONT_SLANT_NORMAL,
                               CAIRO_FONT_WEIGHT_NORMAL);
        cairo_set_font_size(cr, 14);
        cairo_move_to(cr, width / 2 - 80, height / 2);
        cairo_show_text(cr, "No data to display");
        pthread_mutex_unlock(&ctx.graph_mutex);
        return FALSE;
    }

    /* Calculate graph area that we can use*/
    int graph_width = width - margin_left - margin_right;
    int graph_height = height - margin_top - margin_bottom;
    /* Check if it is correct */
    if (graph_width <= 0 || graph_height <= 0)
    {
        pthread_mutex_unlock(&ctx.graph_mutex);
        return FALSE;
    }

    /* Get min and max val for Graphs */
    double min_val1 = ctx.graph_data[0].min_val;
    double max_val1 = ctx.graph_data[0].max_val;
    double min_val2 = ctx.graph_data[1].min_val;
    double max_val2 = ctx.graph_data[1].max_val;

    /* Calculate the range of Y-axis for Graph1*/
    double range1 = (max_val1 > min_val1) ? (max_val1 - min_val1) : 1;
    /* Add padding to Y-axis*/
    uint32_t pad1 = (range1 / 10) + 1;
    /* Calculate new min-max with paddings*/
    min_val1 = (min_val1 > pad1) ? (min_val1 - pad1) : 0;
    max_val1 += pad1;

    /* Calculate the range of Y-axis for Graph2*/
    uint32_t range2 = (max_val2 > min_val2) ? (max_val2 - min_val2) : 1;
    /* Add padding to Y-axis*/
    uint32_t pad2 = (range2 / 10) + 1;
    /* Calculate new min-max with paddings*/
    min_val2 = (min_val2 > pad2) ? (min_val2 - pad2) : 0;
    max_val2 += pad2;

    /* Sliding timing window implementation*/
    uint64_t now = latestSampledTime();

    /* Set min-time */
    uint64_t min_time = (now > SLIDING_WINDOW_US)
                        ? (now - SLIDING_WINDOW_US)
                        : 0;

    uint64_t time_range = SLIDING_WINDOW_US;

    /* Settings of graph and grid lines */
    plotSettings(cr , graph_height , graph_width);

    /* Configuration for Graph-1 */
    plotGraph1 (cr, min_time , time_range , graph_width ,
    					graph_height ,min_val1 , max_val1);


    plotGraph2 (cr , min_time, time_range ,graph_width,
    					graph_height , min_val2, max_val2);

	/* Labeling to Y-Axis */
    //plotLabelYaxis (cr , graph_height , min_val1 , max_val1);
    plotDualYAxis (cr , graph_height ,graph_width ,  min_val1 , max_val1
    				, min_val2 , max_val2);

	/* Labeling to X-Axis */
    plotLabelXaxis (cr , graph_width , graph_height , min_time , time_range);

    /* Axis titles */
    plotTitles(cr , width , height);

    pthread_mutex_unlock(&ctx.graph_mutex);
    return FALSE;
}

\end{lstlisting}



\subsection{GUI Related Requirement Analysis}
This section analyzes the graphical user interface (GUI) related requirements of the 
project and explains how the proposed GUI design and implementation satisfy these 
requirements. Each requirement is evaluated individually, with emphasis on usability, 
layout, user interaction, and responsiveness. The corresponding design choices, interface 
components, and implementation details are discussed. The objective of this analysis is to 
demonstrate that the developed GUI meets the functional, usability, performance, and 
reliability requirements defined at the beginning of the project, while providing a clear 
and intuitive user experience.

\subsubsection{Functional Requirements}
\label{subsubsection:funcreq}

\begin{tcolorbox}[
    title=REQ-9,
    colback=gray!5,
    colframe=gray,
    fonttitle=\bfseries
]
The Control/Monitor program (GUI) shall configure which sensors should acquire data.
\end{tcolorbox}

\paragraph{REQ-9 Analysis and GUI Mapping}\

The GUI provides frequency configuration controls that allow the user to set the 
sampling frequency for each sensor individually. By setting a sensor’s frequency to 0 Hz, 
the corresponding sensor is disabled and stops acquiring data. This requirement is 
satisfied through the implementation of per-sensor frequency configuration within the GUI. 
The functionality can be verified both on real hardware and within a test environment, 
confirming that sensors correctly start or stop data acquisition based on the configured 
frequency.

\begin{tcolorbox}[
    title=REQ-10,
    colback=gray!5,
    colframe=gray,
    fonttitle=\bfseries
]
The Control/Monitor program (GUI) shall configure a sampling rate for each sensor.
\end{tcolorbox}

\paragraph{REQ-10 Analysis and GUI Mapping}\

The GUI allows the user to configure the sampling rate for each sensor by adjusting its 
frequency in the interface. The selected frequency is applied by sending a command to the 
server side, where the command parameters specify which sensor to configure and the 
desired frequency. The server parses these commands through its API and executes the 
corresponding action for the specified sensor. This requirement is verified through
testing on real hardware as well as within a controlled test environment, and further 
details can be investigated by referring to the software development report.


\begin{tcolorbox}[
    title=REQ-11,
    colback=gray!5,
    colframe=gray,
    fonttitle=\bfseries
]

The Control/Monitor program (GUI) shall display live sensor data in text form.
\end{tcolorbox}

\paragraph{REQ-11 Analysis and GUI Mapping}\
The GUI displays live sensor data in text form, positioned on the left side of the 
interface for clear visibility. This functionality is implemented in the 
\text{update\_sensor\_display()} API which can be seen in Listing \ref{lst:update_sensor_display}, which updates the textual display before adding the sensor 
data to the graphical display. This requirement is verified through testing on real 
hardware as well as within a controlled test environment, ensuring that sensor data is 
accurately and consistently displayed in real time.



\begin{tcolorbox}[
    title=REQ-12,
    colback=gray!5,
    colframe=gray,
    fonttitle=\bfseries
]
The Control/Monitor program (GUI) shall display up to 2 signals in a live diagram
\end{tcolorbox}

\paragraph{REQ-12 Analysis and GUI Mapping}\

The GUI displays up to two signals simultaneously in a live graph within the monitor. 
This functionality is implemented in the \text{on\_draw\_graph()} API which can be seen in Listing \ref{lst:draw}. To improve readability and 
visual clarity, the signal values are normalized along the y-axis. This requirement is 
verified through testing on real hardware as well as within a controlled test environment, 
ensuring that up to two signals are accurately displayed in real time.





\subsubsection{Design Requirements}

\begin{tcolorbox}[
    title=REQ-15,
    colback=gray!5,
    colframe=gray,
    fonttitle=\bfseries
]
The GUI shall use GTK C code.
\end{tcolorbox}

\paragraph{REQ-15 Analysis and GUI Mapping}\

The GUI is implemented using the GTK library with the C programming language to ensure 
high performance and efficiency. This implementation choice is discussed in the 
Introduction section, and further details can be found by referring to the relevant 
API documentation presented in this report.


\subsubsection{Test Requirements}

\begin{tcolorbox}[
    title=REQ-26,
    colback=gray!5,
    colframe=gray,
    fonttitle=\bfseries
]
The functional requirents shall be tested by suitable test setups
\end{tcolorbox}

\paragraph{REQ-26 Analysis and GUI Mapping}\

The GUI is tested using multiple test setups to verify all functional attributes of the 
interface. As described in \cref{subsection:asd}, the server is implemented within a Linux-based 
test environment, allowing comprehensive testing of the GUI under various operating 
conditions. All GUI functionalities are verified across different test cases, and the 
corresponding test results are presented in the referenced section.


\subsection{Test and Verification of GUI}
\label{subsection:asd}
This section describes the test and verification process used to validate the 
functionality of the GUI requirements. A Linux-based Ethernet connection is used to 
establish a server on the host machine. Subsequently, test files generating mock 
sensor signals are transmitted to the GUI via the Ethernet interface. This setup enables 
systematic testing and verification of the GUI functionalities under controlled 
conditions. Although it is not feasible to present all possible test cases in this 
section, the most critical and representative cases are discussed.

\subsubsection{Verification of Requirement REQ-11}
First, the overall appearance of the GUI is presented. At this stage, sensor data are 
transmitted from the test environment, and two live graphs are displayed within the GUI.


\begin{figure}[H]
    \centering
    \includegraphics[width=0.8\linewidth]{gui.png}
    \caption{GUI}
    \label{fig:gui}
\end{figure}

First, \cref{fig:gui} is explained. As shown on the left side of the interface, all 
sensor data are displayed in live textual form. Based on this observation, REQ-11 in \cref{subsubsection:funcreq} is verified.

\subsubsection{Verification of Requirement REQ-10}
Secondly, as shown by the sensor frequency configuration row in the upper part of the GUI, 
the sampling frequencies of the sensors can be configured. To verify this functionality, 
the temperature sensor frequency is modified. In Figure \ref{fig:gui}, the temperature sensor operates 
at 500 Hz; this frequency is changed to 10 Hz to examine the resulting signal behavior.

\begin{figure}[H]
    \centering
    \includegraphics[width=0.8\linewidth]{freqchang.png}
    \caption{Frequency Settings Effect to Signal}
    \label{fig:frecch}
\end{figure}

As shown in Figure \ref{fig:frecch} above, the signal quality of the temperature signal is significantly 
reduced. This observation demonstrates that the sampling frequency configuration 
functions as intended; therefore, Requirement REQ-10 is met, as described in Section \ref{subsubsection:funcreq}.

\subsubsection{Verification of Requirement REQ-9}
In addition, sensor sampling can be stopped by setting the corresponding frequency to 0 Hz. 
This functionality corresponds to Requirement REQ-9 described in Section \ref{subsubsection:funcreq}. As shown in 
Figure \ref{fig:zerofreq}, when the ADC0 sensor frequency is set to 0 Hz, the sensor stops sampling and no 
data are displayed. Based on this observation, Requirement REQ-9 is met.

\begin{figure}[H]
    \centering
    \includegraphics[width=0.8\linewidth]{zerfreq.png}
    \caption{Zero Frequency Settings Effect to Signal}
    \label{fig:zerofreq}
\end{figure}

\subsubsection{Verification of Requirement REQ-12}
Requirement REQ-12, which specifies the display of up to two signals in a live diagram, 
is demonstrated by the figures shown above. Additionally, the user can switch between 
different signals by selecting the desired sensors in the graph control section of the GUI. 
Based on these observations, Requirement REQ-12, as defined in Section \ref{subsubsection:funcreq}, is met.

\begin{figure}[H]
    \centering
    \includegraphics[width=0.8\linewidth]{graphcontrol.png}
    \caption{Graph Control for Live Signals}
    \label{fig:graph}
\end{figure}

\subsubsection{Verification of Requirement REQ-15}

This requirement is discussed in the Introduction in section \ref{subsection:intro} and illustrated in Listings \ref{lst:update_sensor_display}
\ref{lst:adddata} and \ref{lst:draw}, 
demonstrating how REQ-15 has been met.

\subsubsection{Verification of Requirement REQ-26}

All of the requirements described above have been verified using the test setups. 
This demonstrates that the functional requirements have been thoroughly tested, 
confirming that Requirement REQ-26 in Section \ref{subsubsection:funcreq} is met.

\subsection{Results}

The GUI was thoroughly tested and verified to ensure all functional requirements are met. 
Sensor data are correctly displayed in both textual and graphical forms, with live updates 
reflecting real-time signals. Users can configure sensor sampling frequencies, including 
stopping sensors by setting their frequency to 0 Hz, and switch between signals on the 
live graph as intended. The implementation using the GTK library in C performs efficiently 
and reliably. Test results from real hardware and controlled environments confirm that 
the GUI functions as designed. Overall, the system meets the expected usability, 
performance, and reliability standards.
